\uuid{kJ5c}
\exo7id{702}
\titre{exo7 702}
\auteur{bodin}
\organisation{exo7}
\datecreate{1998-09-01}
\isIndication{false}
\isCorrection{false}
\chapitre{Dérivabilité des fonctions réelles}
\sousChapitre{Calculs}
\cpge{mpsi-mp2i.derivabilite.a, pcsi.derivabilite.a, ptsi.derivabilite.a, mpsi-mp2i.polynomes-fractions-rationnelles.c, pcsi.polynomes.c, ptsi.polynomes.c, bcpst1.derivees-primitives-integrales.a, bcpst1.derivation.c}
\module{Analyse}
\niveau{L1}
\difficulte{}

\contenu{
\question{
Calculer les d\'eriv\'ees d'ordre $n$ des fonctions :
$$f(x) = \frac{2x-5}{(x-2)^2(x+1)(x-3)}\qquad g(x) = \ln(1+x).$$
}
}
